% Imported from LaTeX_and_pdfs/AGAT_dump.tex; see provenance.json.
\part{Conformal Field Theory; Virasoro Algebra}
  

cf. Schottenloher (2008)  \cite{agat:Scho2008}

\section{Conformal Transformations}

\subsection{Semi-Riemannian manifolds (review and (key) examples)}

cf. pp. 7, Ch. 1 "Conformal Transformations and Conformal Killing Fields." Schottenloher (2008)  \cite{agat:Scho2008}

Semi-Riemannian manifold is a pair $(M, g)$ s.t. \\
smooth manifold $M$, $\text{dim}{M} = n$, \\
smooth tensor field $g$ s.t. $g: a\in M \mapsto \Omega^2(T_aM)$, i.e. 
$\forall \, a \in M$, $g$ assigns $a$ a nonnegative and symmetric bilinear form on tangent space $T_aM$. 

In local coordinates, $x^1 \dots x^n $ of manifold $M$,  \\
given chart $\phi : U \to V$, open subset $U \subseteq M$, open subset $V \subseteq \mathbb{R}^n$, 
\[
\phi(a) = (x^1(a) \dots x^n(a)), \, a\in M
\]
Bilinear form $g_a$ on $T_aM$, written
\[
g_a(X,Y) = g_{\mu \nu}(a) X^{\mu} Y^{\nu}
\]

Tangent vectors $X = X^{\mu} \partial_{\mu}$, $Y = Y^{\nu} \partial_{\nu} \in T_aM$ basis $\partial_{\mu} := \frac{ \partial }{ \partial x^{\mu}}$, $\mu = 1 \dots n$ of tangent space $T_aM$, induced by chart $\phi$. 

By assumption, matrix
\[
g_{\mu \nu}(a)
\]
Nondegenerate and symmetric, $\forall \, a \in U$, i.e. 
\[
\text{det}(g_{\mu \nu }(a)) \neq 0 , \qquad \, (g_{\mu \nu}(a))^T = (g_{\mu \nu}(a))
\]

Differentiating of $g_a$ implies matrix $g_{\mu \nu}(a)$ depends differentiably on $a$.

That means that in its dependence on local coordinates $x^j$, coefficient $g_{\mu \nu} = g_{\mu \nu}(x)$ are smooth functions.

In general, $g_{\mu \nu}X^{\mu} X^{\nu} >0 $ doesn't hold $\forall \, X \neq 0$, i.e. $g_{\mu \nu}(a)$ not required to be positive-definite.

2 important subcases: \footnote{\url{https://doc.sagemath.org/html/en/reference/manifolds/sage/manifolds/differentiable/pseudo_riemannian.html}}

Riemannian manifold: metric $g$ positive definite, signature $n= \text{dim}M$. 

Lorentz manifold specified as semi-Riemannian manifold with $(p,q)= (n-1,1)$ or $(p,q) = (1, n-1)$. 

Metric $g$ has signature $n-2$ (positive convention) or $2-n$ (negative convention).

\subsubsection{Examples (of Riemannian manifolds for Conformal Field Theory)}

$\mathbb{R}^{p,q} = (\mathbb{R}^{p,q}, g^{p,q})$, $p,q \in \mathbb{N}$, where
\[
g^{p,q}(X,Y) := \sum_{i=1}^p X^i Y^i - \sum_{i = p+1}^{p+q} X^i Y^i
\]

Hence
\[
(g_{\mu \nu}) = \left( \begin{matrix} 1_p & \\ & -1_q \end{matrix} \right) = \text{diag}(1\dots 1, -1, \dots  -1)
\]

$\mathbb{R}^{1,3} = \mathbb{R}^{3,1}$, usual Minkowski space. \\
$\mathbb{R}^{1,1}$, 2 -dim. Minkowski space (Minkowski plane). \\
$\mathbb{R}^{2,0}$, Euclidean plane. \\
$\mathbb{S}^2 \subset \mathbb{R}^{3,0}$, compactification of $\mathbb{R}^{2,0}$, structure of Riemannian manifold on 2-sphere $\mathbb{S}^2$ induced by inclusion in $\mathbb{R}^{2,0}$ \\
$\mathbb{S} \times \mathbb{S} \subset \mathbb{R}^{2,2}$, compactification of $\mathbb{R}^{1,1}$. More precisely, \\
$\mathbb{S} \times \mathbb{S} \subset \mathbb{R}^{2,0} \times \mathbb{R}^{0,2} \simeq \mathbb{R}^{2,2}$ where structure of semi-Riemannian manifold on $\mathbb{S} \times \mathbb{S}$ induced by inclusion into $\mathbb{R}^{2,2}$. 

$\mathbb{S}^p \times \mathbb{S}^q \subset \mathbb{R}^{p+1, 0} \times \mathbb{R}^{0, q+1} \simeq \mathbb{R}^{p+1,q+1}$ with $p$-sphere $\mathbb{S}^p = \lbrace X \in \mathbb{R}^{p+1} : g^{p+1, 0}(X,X) = 1\rbrace \subset \mathbb{R}^{p+1,0}$, $q$-sphere $\mathbb{S}^q \subset \mathbb{R}^{0, q+1}$ \\
yields a compactification of $\mathbb{R}^{p,q}$ for $p,q \geq 1$

Compact semi-Riemannian manifold denoted by $\mathbb{S}^{p,q}$, for $p, q\geq 0$.

Quadrics $N^{p,q}$ (of Sec. 2.1) are locally isomorphic to $\mathbb{S}^{p,q}$ from point of view of conformal geometry.

For the "\emph{negative convention}": 

\[
g^{p,q}(X,Y) = - \sum_{i=0}^{p-1} X^i Y^i + \sum_{i=p}^{p+q} X^i Y^i
\]
\[
(g_{\mu \nu}) = \left( \begin{matrix} -1_p & \\ & 1_q \end{matrix} \right) = \text{diag}(-1 , \dots - 1, 1 \dots 1)
\]

$\mathbb{R}^{1,3} $, Minkowski space. \\
$\mathbb{R}^{1,1}$, 2 -dim. Minkowski space. \\
$\mathbb{R}^{0,2}$, Euclidean plane. \\
$\mathbb{S}^2 \subset \mathbb{R}^{0,3}$, compactification of $\mathbb{R}^{0,2}$ \\
$\mathbb{S} \times \mathbb{S} \subset \mathbb{R}^{0,2} \times \mathbb{R}^{2,0} \simeq \mathbb{R}^{2,2}$ 

$\mathbb{S}^p \times \mathbb{S}^q \subset \mathbb{R}^{0, p+1} \times \mathbb{R}^{q+1, 0} \simeq \mathbb{R}^{p+1,q+1}$ with $p$-sphere $\mathbb{S}^p = \lbrace X \in \mathbb{R}^{p+1} : g^{ 0, p+1}(X,X) = 1\rbrace \subset \mathbb{R}^{0, p+1}$, $q$-sphere $\mathbb{S}^q \subset \mathbb{R}^{q+1, 0}$ \\
yields a compactification of $\mathbb{R}^{p,q}$ 




\begin{definition}[Conformal transformation or conformal map]
Let 2 semi-Riemannian manifolds$(M,g)$, $(M',g')$, $\text{dim}M = \text{dim}M'$, let open $U\subset M$, open $V\subset M'$.  

\textbf{conformal transformation} or \textbf{conformal map} is a smooth $\varphi : U\to V$ of maximal rank, if $\exists \, $ smooth $\Omega : U\to \mathbb{R}^+$ s.t. 
\begin{equation}
	\varphi^* g' = \Omega^2 g
\end{equation}
where $\varphi* g'(X,Y) :=g'(T\varphi(X), T\varphi(Y))$ and $T\varphi: TU \to TV$ denote tangent map (derivative) of $\varphi$.  

$\Omega \equiv $ \emph{conformal factor} of $\varphi$.  
\end{definition}

Locally, $y^i = \varphi^i(x)$, \[
\begin{gathered}
	\frac{ \partial \varphi^i }{ \partial x^j} = \frac{ \partial y^i }{ \partial x^j}
\end{gathered}	
\]
Then
\[
X = X^k \frac{ \partial }{ \partial x^k} = X^k \frac{ \partial y^i }{ \partial x^k} \frac{ \partial }{ \partial y^i } = X^k \frac{ \partial \varphi^i }{ \partial x^k} \frac{ \partial }{ \partial y^k} \in TM
\]
and so 
\[
\begin{gathered}
	\varphi^* g'(X,Y) = g'(T\varphi(X), T\varphi(Y)) = (g')_{ij} X^k \frac{ \partial y^i }{ \partial x^k} Y^l \frac{ \partial y^j }{ \partial x^l} = (g')_{ij} X^k \frac{ \partial \varphi^i }{ \partial x^k } Y^l \frac{ \partial y^j}{ \partial x^l} \\ 
\Longrightarrow (\varphi^* g')_{kl} = (g')_{ij} \frac{ \partial y^i }{ \partial x^k } \frac{ \partial y^j}{ \partial x^l } \\ 
\Longrightarrow (\varphi^* g')_{kl} = (g')_{ij} \frac{ \partial \varphi^i }{ \partial x^k} \frac{ \partial \varphi^j}{ \partial x^l} = \Omega^2 g_{kl}
\end{gathered}
\]




\begin{definition}
	\textbf{extension} of $G$ by group $A$ is (given by) an \emph{exact sequence} of group homomorphisms.  
	\begin{tikzpicture}
	\matrix(m)[matrix of math nodes, row sep=4em, column sep=4em]
	{
		1   &  A & E & G & 1 \\
	};
	\path[->]
	(m-1-1) edge node[auto]{$$} (m-1-2)
	(m-1-2) edge node[auto]{$i$} (m-1-3)
	(m-1-3) edge node[auto]{$\pi$} (m-1-4)
	(m-1-4) edge node[auto]{$$} (m-1-5)	
	;
	\end{tikzpicture} 
\end{definition}
cf. Def. 3.1 of Schottenloher (2008)  \cite{agat:Scho2008}.  

Recall that an exact sequence, if $\begin{aligned} & \quad \\
	& \text{im}(1\to A) = \text{ker}(i) \\
	& \text{im}(i) = \text{ker}(\pi) \\
	& \text{im}(\pi) = \text{ker}(G\to 1)
\end{aligned}$

By Thm., 
$1\to A \xrightarrow{i} E$ exact so $i$ injective.  \\
$E\xrightarrow{\pi} G \to 1$ exact so $\pi$ surjective.  

Extension is called \textbf{central} if $A$ abelian and image $\text{im}i$ is in center of $E$, i.e. 
$a\in A, b\in E \Longrightarrow i(a) b = bi(a)$.  

\subsubsection{Examples of extensions of $G$, and central extensions of $G$ (which has a particular $E$)}
\begin{itemize}
\item e.g. central extension has form 
\[
\begin{tikzpicture}
\matrix(m)[matrix of math nodes, row sep=4em, column sep=4em]
{
	1   &  A & A\times G & G & 1 \\
};
\path[->]
(m-1-1) edge node[auto]{$$} (m-1-2)
(m-1-2) edge node[auto]{$i$} (m-1-3)
(m-1-3) edge node[auto]{$\text{pr}_2$} (m-1-4)
(m-1-4) edge node[auto]{$$} (m-1-5)	
;
\end{tikzpicture} 
\]
where  \\
$\begin{aligned}
 i:A & \to A\times G \\
a & \mapsto (a,1) \end{aligned}$ 

\[
\begin{gathered}
i(a)(a',g) = (a,1)(a',g) = (aa',g) = \\
= (a'a,g\cdot 1) = (a',g)(a,1) = (a',g)i(a)
\end{gathered}
\]

Notice that what the \emph{exactness} property of an exact sequence does:
\[
\text{pr}_2i(a) = \text{pr}_2(a,1) = 1
\]

\item e.g. of a \text{nontrivial central extension} is exact sequence 
\begin{equation}
\begin{tikzpicture}
\matrix(m)[matrix of math nodes, row sep=4em, column sep=4em]
{
	1   &  \mathbb{Z}/k\mathbb{Z} & E\times U(1) & U(1) & 1 \\
};
\path[->]
(m-1-1) edge node[auto]{$$} (m-1-2)
(m-1-2) edge node[auto]{$$} (m-1-3)
(m-1-3) edge node[auto]{$\pi$} (m-1-4)
(m-1-4) edge node[auto]{$$} (m-1-5)	
;
\end{tikzpicture} 
\end{equation}
with $\pi(z) = z^k$ \, $\forall \, k \in \mathbb{N}$, $k\geq 2$, since $E=U(1)$ and $\mathbb{Z}/k\mathbb{Z}$ are not isomorphic.  

Also, homomorphism $\tau:U(1) \to E$ with $\pi \circ \tau = 1_{U(1)}$, doesn't exist, since there's no global $k$th root.  

EY : 20170926 It's that in integer division of the argument in a complex number $z\in U(1)$, and exponent multiplication by $k$, you go from 1 to many and many to 1, depending upon the "branch" you're mapping to for complex numbers.  

For $[n] \in \mathbb{Z}/k\mathbb{Z}$, 
\[
[n]\xmapsto[]{ i } \exp{ \left( \frac{ [n] }{ k} 2\pi i \right) }
\]
and so 
\[
\text{ker}\pi = \lbrace z | \pi(z) = 1 \rbrace \text{ so that } \text{ker}\pi = \lbrace z = \exp{ \left( \frac{i 2\pi n}{k} \right) } \rbrace
\]

\item e.g. \emph{Semidirect products}.  

group $G$ acting on another group $H$, by homomorphism   
\[
\tau : G \to \text{Aut}(H)
\]






\begin{definition}[semi-direct product]\label{agat:Def:Semidirectprod}
\textbf{semidirect product} group $G \ltimes H$ is set $H\times G$, with multiplication 
\[
(x,g) \cdot (x',g') := (x\tau(g)(x'), gg') \qquad \, \forall \, (x,g), (x',g') \in H\times G
\]
\end{definition}
\begin{equation}
\begin{tikzpicture}
\matrix(m)[matrix of math nodes, row sep=4em, column sep=4em]
{
	1   &  H &  G \ltimes H & G & 1 \\
};
\path[->]
(m-1-1) edge node[auto]{$$} (m-1-2)
(m-1-2) edge node[auto]{$i$} (m-1-3)
(m-1-3) edge node[auto]{$\pi$} (m-1-4)
(m-1-4) edge node[auto]{$$} (m-1-5)	
;
\end{tikzpicture} 
\end{equation}
with 
\begin{equation}
	\begin{aligned}
	& i : H\to G \ltimes H \\ 
	& i(x) = (x,1)
\end{aligned}
\end{equation}
$i$ group homomorphism, since 
\[
i(x_1x_2) = (x_1x_2 ,1) = (x_1 \tau(1)x_2 , 1) = (x_1,1)\cdot (x_2,1) = i(x_1)i(x_2)
\]
\begin{equation}
	\begin{aligned}
	& \pi : G \ltimes H \to G\\ 
	& \pi(x,g) = g 
\end{aligned}
\end{equation}
cf. \url{http://sierra.nmsu.edu/morandi/oldwebpages/math683fall2002/GroupExtensions.pdf}

Observe that 
\[
\pi i(x) = \pi(x,1)=1 \text{ so } \text{ker}\pi = \text{im}i
\]

\begin{definition}[Semi-direct product (2); with direct product]\label{agat:Def:Semidirectprod2}
	\textbf{direct product} $G=HK$ if \\
$H,K$ subgroups of group $G$, s.t. \\
\begin{itemize}
\item $H$ and $K$ are normal in $G$ ($gkg^{-1} \in K$ \, $\forall \, g\in G$, $\forall \, k \in K$)  
\item $H\cap K = \lbrace 1 \rbrace$  
\item $HK=G$.  
\end{itemize} 

\textbf{semi-direct product}.  Relax the 1st condition (of direct products) so $H$ still normal in $G$, but $K$ need not be.  \\
\begin{itemize}
\item $H$ normal in $G$ ($ghg^{-1} \in H, \, \forall \, g , \, \forall \, h \in H$)  
\item $H\cap K = \lbrace 1 \rbrace$ 
\item $HK = G$
\end{itemize}
\end{definition}

\emph{Connection between} Def. \ref{agat:Def:Semidirectprod} and Def. \ref{agat:Def:Semidirectprod2} for the \emph{semidirect product}:  Consider $\tau : G \to \text{Aut}(H)$.  

Consider $G \ltimes H$ - what is the identity $1_{G\ltimes H} \equiv (1_H,1_G)$ of this group?  
\[
(x,g)\cdot (1_H,1_G) = (x\tau(g) 1_H, g1_G) = (x\tau(g)1_H, g) \Longrightarrow 1_H = \tau(g^{-1})1, \, 1_G=1
\]
and so the inverse, $\forall (x,g) \in G\ltimes H$, $(x,g)^{-1} \equiv ((x^{-1}), (g^{-1}))$, 
\[
\begin{gathered}
	(x,g) (x,g)^{-1} = (x\tau(g) (x^{-1}),g(g^{-1}) ) = (x\tau(g)(x^{-1}), 1) \text{ (if $(g^{-1}) = g^{-1}$) }
\end{gathered}
\]
Moving along, 
\[
\begin{gathered}
	x\tau(g) (x^{-1}) = \tau(g^{-1})1 \\
\Longrightarrow (x^{-1}) = \tau(g^{-1}) x^{-1}\tau(g^{-1})1
\end{gathered}
\]

Checking out the $H$ being a normal subgroup of $G\ltimes H$ condition, i.e. $H \lhd G$, 
\[
\begin{gathered}
	(x,g)(h,1)(\tau(g^{-1})x^{-1} \tau(g^{-1}) , g^{-1}) = (x\tau(g)h, g)(\tau(g^{-1})x^{-1}\tau(g^{-1}) , g^{-1}) = \\
= (x\tau(g)h \tau(g)\tau(g^{-1}) x^{-1}\tau(g^{-1}), 1) = (x\tau(g) hx^{-1}\tau(g^{-1}), 1) 
\end{gathered}
\]
$\Longrightarrow$ $H$ normal subgroup of $G\ltimes H \equiv H \lhd (G\ltimes H)$.  

\href{http://www.math.columbia.edu/~bayer/S09/ModernAlgebra/semidirect.pdf}{Notes on Semidirect products}
\item extension 
\begin{equation}
\begin{tikzpicture}
\matrix(m)[matrix of math nodes, row sep=4em, column sep=4em]
{
	1   &  SL(n,\mathbb{R}) &  GL(n,\mathbb{R}) & \mathbb{R}^* & 1 \\
};
\path[->]
(m-1-1) edge node[auto]{$$} (m-1-2)
(m-1-2) edge node[auto]{$i$} (m-1-3)
(m-1-3) edge node[auto]{$\text{det}$} (m-1-4)
(m-1-4) edge node[auto]{$$} (m-1-5)	
;
\end{tikzpicture} 
\end{equation}

with \\
$GL(n,\mathbb{R}) \equiv Gl_n(\mathbb{R}) = \lbrace A | A \in \text{Mat}_{\mathbb{R}}(n,n); \text{det}A \neq 0\rbrace$ \\
$\text{det}:GL(n,\mathbb{R}) \to \mathbb{R}^* \equiv \mathbb{R} \backslash \lbrace 0 \rbrace $, $\text{det}$ surjective homomorphism \\
$SL(n,\mathbb{R}) \equiv Sl_n(\mathbb{R}) = \lbrace A | A \in \text{Mat}_{\mathbb{R}}(n,n); \, \text{det}A=1\rbrace$ \\

Note that $\text{ker}(\text{det}) = SL(n,\mathbb{R})$.  

Now 
\[
\mathbb{R}^* \simeq \lbrace a 1_n | a \in \mathbb{R}^* \rbrace
\]
and $\text{det}(a1_n) = a^n$.  

If $n$ odd, and $\text{det}(a1_n)=a^n =1$, then $a=1$.  If $n$ even, $a=\lbrace -1, 1 \rbrace$.  

By the second definition of a semi-direct product, Def. \ref{agat:Def:Semidirectprod2}, it's required that $SL(n,\mathbb{R}) \cap \mathbb{R}^* = 1$ (i.e. the intersection is only the identity).  This will only be the case if $n$ odd.  

cf. \url{http://sierra.nmsu.edu/morandi/oldwebpages/math683fall2002/GroupExtensions.pdf}

\end{itemize}

